\chapter{Acceptance criteria, registered in advance}
\label{chap:criteria}

This chapter is written before the session. Its thresholds are fixed here so
that the result can be judged against a standard that was not chosen with the
result already in view. Where a threshold is derived from physics, the
derivation is given. Where it is a project decision, it is labelled as one and
is open to revision, but only by amending this chapter in a dated revision,
not silently at analysis time.

\section{Why this chapter is not optional}

A low residual on a degenerate set is not evidence of a good calibration; it
is evidence that the ambiguity was absorbed somewhere. The project has already
been caught by exactly this: the best archived session yields a root-mean-square
reprojection error of 1.50~px, which the existing code classified as merely
high rather than catastrophic, while the same fit gives
$f_x = 2651.79 \pm 22.42$~px, a 0.85\% uncertainty, and a distortion
coefficient $k_3 = 3.129 \pm 3.714$ whose standard deviation exceeds its own
value \cite{linku_calibration_method}. A number that looks acceptable while
the parameters behind it are undetermined is the specific failure this chapter
is designed to prevent repeating at the level of the whole experiment.

\section{A1. Calibration}

\begin{table}[htbp]
\centering
\caption[Calibration acceptance criteria]{Acceptance criteria for the
calibration session that precedes the experiment. All are evaluated by the
guided tool at the end of the session, before any experimental capture is
taken.}
\label{tab:crit-calib}
\begin{tabularx}{\textwidth}{@{}l X l@{}}
\toprule
\textbf{Id} & \textbf{Criterion} & \textbf{Threshold} \\
\midrule
A1.1 & Accepted views covering each of the four tilt bins
       (0--10\textdegree, 10--30\textdegree, 30--45\textdegree,
       45--70\textdegree) & $\geq 3$ per bin \\
A1.2 & Relative standard deviation of $f_x$ and of $f_y$ & $\leq 0.3\%$ \\
A1.3 & Standard deviation of each distortion coefficient, relative to the
       magnitude of that coefficient & $< 1$ \\
A1.4 & Offset of the principal point from the image centre, as a fraction of
       image width & $\leq 5\%$ \\
A1.5 & Capture resolution of the calibration frames & $4056 \times 3040$ \\
A1.6 & Cameras not moved, and focus and aperture rings not touched, between
       calibration and the last experimental capture & binary \\
\bottomrule
\end{tabularx}
\end{table}

\textbf{Derivation of A1.2.} The focal length in pixels sets the scale of
every reconstructed distance directly: a relative error $\delta$ in $f_x$
produces a relative error $\delta$ in every depth. At the nominal working
distance of 3.5~m, holding the scale error from this source below 1~cm
requires $\delta \leq 0.29\%$, which rounds to the registered 0.3\%. The
archived best case is 0.85\%, about three times this.

This is the one threshold in the document that may prove unreachable on this
rig. If it is, the correct response is to record the value achieved and state
that the scale error is correspondingly larger, not to relax the number
afterwards. The figure stands as the requirement the measurement has to be
judged against, whether or not it is met.

\textbf{Derivation of A1.3.} A parameter whose standard deviation exceeds its
own magnitude is not determined by the data; the fit is free to set it to zero
without materially changing the residual. This is a binary sanity check, not a
tuned threshold, and the archived fit fails it on $k_3$.

\textbf{Derivation of A1.4.} A principal point far from the image centre is
the signature of the $[t_x, t_y]$ against $[c_x, c_y]$ degeneracy, in which a
board translation and a principal-point shift become indistinguishable. The
active calibration shows $c_x = 694.8$ on a 1920-pixel-wide image where it
should sit near 960, an offset of 265~px or 13.8\% of the width. The 5\%
threshold is a project decision chosen to sit comfortably below the known
failure and comfortably above the few-percent decentring a real lens
assembly can legitimately show.

\textbf{On A1.5.} The relation $f_x = f / p$ means a calibration is tied to
the effective pixel size of the mode it was captured in, and two of the
IMX477 modes crop as well as bin, so the result cannot simply be rescaled
\cite{raspberrypi_camera_docs}. Every archived session used
$1920 \times 1440$, never the full sensor. This criterion exists because the
single most common way to waste the session is to calibrate in one mode and
capture in another.

\section{A2. Illumination}

Each requirement is taken from the dark-room specification
\cite{linku_darkroom_illumination}, where its derivation from primary sources
is given. Reproduced here as pass or fail conditions with the measurement that
decides them.

\begin{table}[htbp]
\centering
\caption[Illumination acceptance criteria]{Illumination acceptance criteria
and the measurement that decides each.}
\label{tab:crit-light}
\begin{tabularx}{\textwidth}{@{}l X l@{}}
\toprule
\textbf{Id} & \textbf{Criterion and how it is measured} & \textbf{Threshold} \\
\midrule
A2.1 & Stray light: with the fixture off, at the longest exposure to be used,
       no scene structure discernible in either camera & nothing visible \\
A2.2 & Single direction: one fixture, or several with parallel axes. No fill,
       no bounce, no ambient & binary \\
A2.3 & Collimation: width of the penumbra cast by a straight edge held 1~m
       from a white card at the tether position & $\leq 35$~mm \\
A2.4 & Uniformity over the working volume, including the depth axis, as
       $(\max-\min)/(\max+\min)$ on a lux grid & $\leq 5\%$ \\
A2.5 & Illuminance at the tether position & $\geq 2320$~lx \\
A2.6 & Short-term stability over the session & $\leq 2\%$ \\
A2.7 & Flicker: white card imaged at the shortest exposure in use, each
       camera & no banding \\
A2.8 & Colour temperature & 5600--6000~K \\
\bottomrule
\end{tabularx}
\end{table}

\textbf{On A2.3.} The real Sun subtends 0.533\textdegree{} and would cast a
9.3~mm penumbra at 1~m \cite{prsa2016iau}. The registered 35~mm corresponds to
a source of about 2\textdegree, set just outside the 1.9\textdegree{} of the
ESA Large Space Simulator \cite{esa_lss}. The virtue of this test is that it
needs a ruler rather than an optical bench.

\textbf{On A2.4.} This is the class B uniformity of IEC 60904-9, using that
standard's own definition \cite{iec60904_9_2020}. The depth axis is the one
that fails first and the one most often omitted from a uniformity check. The
procedure in Chapter \ref{chap:procedure} loops back on failure of this
criterion specifically, because moving a fixture to fix uniformity changes
the illuminance while dimming it does not change the uniformity.

\textbf{On A2.5.} Derived from the Sony IMX477 datasheet
\cite{sony_imx477_datasheet} for the fitted 8~mm f/1.6 configuration at the
1\textdegree/s reference rate, as the illuminance at which a white surface
reaches saturation within the maximum exposure that rate allows. It is a
floor, not a target: four stated assumptions in the derivation all push the
real requirement upward, and a tether darker than a white card needs more.

\textbf{On A2.7.} A mains-connected fixture modulates at twice the line
frequency, giving a 10~ms cycle at 50~Hz, the same order as the exposures in
use, and a rolling shutter turns that into banding
\cite{sheinin2018rollingac}. Note that a global-shutter sensor is not immune:
it converts the banding into frame-to-frame brightness variation, which is
equally destructive for a stereo pair that has to be radiometrically
consistent. Tuning the exposure to an exact multiple of the flicker period is
not available as a remedy, because the exposure is not free. The fixture
itself must not flicker.

\section{A3. Pipeline}

These are evaluated after the session, on the captures, by running the
application unmodified.

\begin{table}[htbp]
\centering
\caption[Pipeline acceptance criteria]{Pipeline acceptance criteria. Unless
stated otherwise each is evaluated on every stereo pair selected for analysis,
and the criterion is met if the stated fraction of pairs passes.}
\label{tab:crit-pipeline}
\begin{tabularx}{\textwidth}{@{}l X l@{}}
\toprule
\textbf{Id} & \textbf{Criterion} & \textbf{Threshold} \\
\midrule
A3.1 & Segmentation yields a tether mask that is a single connected component
       in both images & $\geq 80\%$ of pairs \\
A3.2 & Endpoint detection returns exactly two endpoints in both images &
       $\geq 80\%$ of pairs \\
A3.3 & Automatic tracking reaches the far endpoint without operator
       intervention & $\geq 50\%$ of pairs \\
A3.4 & Fraction of path samples surviving the disparity and depth filters of
       \texttt{compute\_wire\_metrics} & $\geq 90\%$ per pair \\
A3.5 & Reconstructed arc length against the tape-measured tether length, on
       pairs where the whole tether is in frame in both cameras &
       within $\pm 10\%$ \\
A3.6 & Straightness, on a pair where the tether was deliberately held taut &
       $\geq 0.97$ \\
A3.7 & At least one complete run producing an exported point cloud and a full
       set of metric scalars & $\geq 1$ \\
\bottomrule
\end{tabularx}
\end{table}

\textbf{On A3.3.} The 50\% threshold is deliberately lower than the others.
Automatic tracking is the stage most exposed to crossings and to breaks in the
mask, and the application offers a hand-drawn path as a first-class
alternative precisely because the automatic route is not expected to be
universally reliable yet. A result below 50\% is informative rather than
disqualifying: it is the measurement that would justify building the learned
branch scorer.

\textbf{On A3.4.} \texttt{compute\_wire\_metrics} discards points with
non-positive disparity or with $Z$ outside 0.05 to 50~m \emph{before}
computing arc length, so a mid-tether segment that gets filtered makes the
reported total length too short with no warning anywhere in the output. The
retained fraction therefore has to be recorded alongside the length, or the
length cannot be interpreted. This criterion exists to force that.

\textbf{On A3.5.} This is the only criterion in the document that compares
against an external measurement. The tether length is measured with a tape
before the session and recorded in Chapter \ref{chap:record}. The
$\pm 10\%$ tolerance is a project decision and is loose on purpose for a first
run: it has to accommodate the focal scale error permitted by A1.2, the
systematic error of the parametric matching assumption, and the under-reporting
caused by A3.4, none of which has ever been measured together on hardware. It
is set to detect a gross failure, not to certify an accuracy. Tightening it is
the job of the next revision, once there is a measured distribution to tighten
it against.

\textbf{On A3.6.} Straightness is reported by the pipeline as the ratio of
the endpoint chord to the traced arc length, so a taut tether should return a
value just below 1. The 0.97 threshold allows roughly 3\% of path wander from
sag, mask noise and tracking jitter combined.

\section{A4. Record}

\begin{table}[htbp]
\centering
\caption[Record acceptance criteria]{Completeness criteria for the session
record. These decide whether the experiment can be replicated and rendered,
not whether it succeeded.}
\label{tab:crit-record}
\begin{tabularx}{\textwidth}{@{}l X@{}}
\toprule
\textbf{Id} & \textbf{Criterion} \\
\midrule
A4.1 & Every field of the record sheet in Chapter \ref{chap:record} filled, or
       explicitly marked not measured with the reason \\
A4.2 & Tether physical pose recorded: approach angle, distance to the camera
       baseline, and attachment points \\
A4.3 & Fixture make, model, modifier, count and ballast or driver mode
       recorded \\
A4.4 & Ring marks checked and photographed at the end of the session, and the
       result recorded \\
A4.5 & Calibration session and all experimental captures committed to the
       repository, with their per-pair metadata \\
\bottomrule
\end{tabularx}
\end{table}

\textbf{On A4.4.} The experiment moves the rig by hand. Moving the focus ring
changes the focal length, and $f_x$ is the focal length in pixels, so 1\% of
drift is 1\% of scale error on every reconstructed distance. It is a
systematic bias and averaging more frames does not remove it. If the marks
have moved by the end of the session, the data is in doubt and A3.5 cannot be
interpreted; this is the reason the check is a criterion and not a note.

\textbf{On A4.5.} The Blender validation deliverable records that for its
chosen reference pair the capture session and exported point cloud were never
pushed to the repository, leaving only screenshots
\cite{linku_blender_validation}. This criterion exists so that it does not
happen twice.

\section{Overall verdict}

The experiment is reported as a \textbf{pass} if A1, A2 and A4 are met in full
and A3 is met except for A3.3 and A3.5. It is reported as a \textbf{qualified
pass} if A1 and A4 are met, A2 fails on no more than one of A2.3, A2.4 or
A2.8, and A3.7 is met: in that case the captures are usable for geometry but
not as the photometric reference for the Blender comparison. It is reported as
a \textbf{fail} if A1 is not met, because then no quantity in metres means
anything, or if A3.7 is not met, because then the pipeline did not complete
even once.
